The generation of synthetic financial time series that faithfully preserve the statistical properties of real market data is a central challenge in quantitative finance and data quality research. High-fidelity synthetic data are essential for stress testing risk models against plausible but unobserved market scenarios, validating portfolio optimization algorithms, and augmenting limited training sets for machine learning pipelines \cite{assefa_generating_2021, jordon_synthetic_2022}. Yet the quality requirements for financial synthetic data are unusually demanding: empirical equity excess growth rates exhibit three well-documented statistical regularities, namely heavy-tailed (leptokurtic) distributions, negligible linear autocorrelation in raw returns, and persistent volatility clustering, that any generative framework must simultaneously reproduce to be considered statistically faithful \cite{mandelbrot_variation_1963, cont_empirical_2001, bauer_jdiq_ts_synthesis_2024}. These domain-specific quality criteria, often called \textit{stylized facts}, serve as the benchmark against which the fidelity of synthetic financial data must be evaluated.

Existing generative approaches address these quality requirements from different angles but each falls short on at least one dimension. GARCH-family models \cite{engle_1982, bollerslev_garch_1986} capture conditional variance dynamics but do not explicitly represent discrete market regimes or sudden jumps in volatility. Stochastic-volatility and jump-diffusion models \cite{heston_stochastic_1993, merton_jump_1976} enrich tail behavior but lack a mechanism for volatility persistence after jumps. Deep generative models, including recurrent neural networks \cite{hochreiter_lstm_1997, fischer_krauss_2018} and generative adversarial networks \cite{takahashi_2019, kwon_can_2024, yoon_timegan_2019}, can learn complex distributions from data but struggle to reproduce temporal dependence structures, particularly volatility clustering, without explicitly encoding the relevant inductive biases \cite{bauer_jdiq_ts_synthesis_2024}. Hidden Markov models (HMMs) and related regime-switching frameworks offer a probabilistic structure for non-stationary market behavior through latent state dynamics \cite{rabiner_introduction_1986}, yet standard HMM specifications fail to generate persistent high-volatility regimes, leading to overly rapid reversion following extreme events and poor temporal quality in the synthetic output \cite{ryden_stylized_1998, bulla_stylized_2006}. Bridging this gap requires a generative framework that preserves both distributional fidelity and temporal structure while remaining computationally scalable.

We address this challenge by developing a hybrid HMM framework that discretizes continuous excess growth rates into quantile-based regimes and augments the resulting Markov process with a jump-duration mechanism that enforces volatility persistence, linking smooth diffusion to discontinuous shocks identified by jump-detection methods \cite{au_yeung_jump_2020}. Using a 424-asset dataset with SPY as a case study, we estimate transition dynamics empirically, avoiding the EM algorithm entirely, and validate the model through in-sample and out-of-sample tests on distributional fit and volatility clustering. We further extend the univariate framework to a correlated multi-asset setting via a Single-Index Model (SIM) factor structure \cite{sharpe_simplified_1963}, enabling synthetic path generation across all 424 assets in a single generative pass.

This paper makes four contributions to synthetic financial data generation and quality assessment. First, we propose augmenting a discrete-state regime-switching model with a two-parameter Poisson jump-duration mechanism $(\epsilon, \lambda)$ that forces the model to dwell in high-volatility tail states for empirically realistic durations, providing a structural inductive bias that preserves the temporal quality (volatility clustering) of the synthetic output, addressing the key limitation identified in prior evaluations of both HMM-based \cite{ryden_stylized_1998} and GAN-based \cite{kwon_can_2024} financial data generators. Second, by partitioning excess growth rates into quantile-defined states via a Laplace CDF, we enable direct frequentist counting of regime transitions, eliminating the computational cost and initialization sensitivity of the Baum-Welch EM algorithm and making the framework scalable to a 424-asset synthetic data pipeline. Third, we conduct a rigorous quality evaluation using three complementary metrics (Kolmogorov-Smirnov and Anderson-Darling pass rates for distributional fidelity, and ACF mean absolute error for temporal structure preservation), all reported with standard errors. Using ten years of SPY data (2014--2024) for training and 249 trading days (full calendar year 2025) for out-of-sample testing, we demonstrate KS and AD pass rates of 97.1\% (SE 0.5) and 90.4\% (SE 0.9) in-sample and 95.0\% (SE 0.7) and 95.4\% (SE 0.7) out-of-sample across 1,000 simulated paths, with volatility clustering preserved in both windows. Finally, we extend the univariate framework to the full 424-asset universe through a Single-Index Model factor projection, enabling scalable correlated synthetic path generation while preserving cross-sectional correlation structure.

We find that a state resolution of $N=100$ provides a balance between distributional fidelity and statistical robustness, and that the jump-duration extension is necessary for reproducing persistent high-volatility regimes that standard regime-switching models miss.

% ── Figure 1: Empirical Motivation ───────────────────────────────────────────
\begin{figure}[htpb]
    \centering
    \includegraphics[width=\textwidth]{Fig1-Empirical-Motivation.pdf}
    \caption{Empirical stylized facts of SPY daily excess growth rates (2014--2024).
    Panel~(a) shows the leptokurtic marginal distribution; the Laplace fit substantially outperforms
    the Gaussian.  Panel~(b) confirms heavy tails via a normal Q-Q plot.  Panels~(c) and~(d)
    contrast the near-zero autocorrelation of raw returns (consistent with the efficient markets
    hypothesis) against the persistent autocorrelation of absolute returns, motivating the
    jump-duration extension.}
    \label{fig:empirical_motivation}
\end{figure}
